% report_facct.tex — FAccT/ACM draft. PRIMARY paper for this project.
%
% Thesis: at every layer of a fairness audit — the metric library, the mapping
% from attribute to documented construct, and the instruments that evaluate
% the audit's own components — the construct actually operationalised differs
% from the construct named, the output is clean either way, and nothing
% records the substitution. The paper makes measurement claims only; law
% appears solely as the provenance of definitions and flags, never
% interpreted (user decision 2026-08-10: no legal stance).
%
% The NeurIPS Datasets & Benchmarks variant is report/report.tex (same project,
% artifact-first framing; shares refs.bib and all numbers). Numbers changed
% here must be propagated there.
%
% Structure (FAccT-standard scaffold, 2026-09 restructure): Introduction /
% Related Work / Methods / Results (the three layers plus measurement
% validity, as subsections) / Discussion / Recommendations / Limitations /
% Conclusion. Section labels (sec:library, sec:statute, ...) were kept
% through the restructure, so cross-references and the companion docs
% (ARGUMENT_SKELETON.md, PAPER_MAP.md in this directory) remain stable.
%
% Every number in this paper traces to a regenerable artifact; appendix G maps
% claim -> artifact -> command, and tests/test_artifact_map.py asserts the
% artifacts exist. Tables under report/generated/ are emitted by scripts and
% must not be edited by hand.
%
% Build:  make paper-facct   (tectonic; it runs BibTeX itself)
%
% For double-blind submission switch to:
%   \documentclass[sigconf,anonymous,review]{acmart}
% `nonacm` is set because this draft has no assigned DOI/ISBN.

\documentclass[sigconf,nonacm]{acmart}

\usepackage{booktabs}
\usepackage{multirow}
\usepackage{verbatim}

\citestyle{acmnumeric}

\begin{document}

\title{The Metric Is Not the Construct}
\subtitle{Construct Validity at Every Layer of a Lending Audit}

\author{Siddhant Jain}
\affiliation{%
  \institution{Simon Fraser University}
  \city{Burnaby}\state{British Columbia}\country{Canada}}

\author{Ayaan Lalani}
\affiliation{%
  \institution{Simon Fraser University}
  \city{Burnaby}\state{British Columbia}\country{Canada}}

\author{Nicholas Vincent}
\affiliation{%
  \institution{Simon Fraser University}
  \city{Burnaby}\state{British Columbia}\country{Canada}}

\renewcommand{\shortauthors}{Jain, Lalani, and Vincent}

\begin{abstract}
The same model, evaluated on the same applicants, can pass or fail a fairness
audit depending on settings that appear nowhere in the audit's output. This
paper identifies those settings across a complete audit of consumer lending
($1{,}000$ credit applications and $10{,}978$ Georgia mortgage applications)
and measures how often each one changes the verdict. Each is an instance of
one underlying event: a construct someone intended (``favorable outcome,''
``the protected class,'' ``refusal'') is operationalised as something a
program can compute, the two differ, and the output looks equally clean
either way. Every wrong verdict we measured erred in the same direction:
toward the clean, reassuring reading.

The settings live at three layers. In the metric library (verified directly
in AIF360~0.6 and Fairlearn~0.13, and scoped to what we verified): one AIF360
constructor defaults \texttt{favorable\_label=1.0}, so on a target named
\texttt{loan\_default} it reports the disadvantaged group as favoured, with
$\mathrm{DI}=1.2072$ where the oriented value is $0.9862$; and neither
library attaches an interval, a sample size, or a record of the reference
group to a disaggregated ratio. In the mapping from attribute to construct: a
banded age attribute reads $0.9791$, no disparity, where the $62{+}$ cut (the
boundary the mortgage data ship a dedicated flag for) reads $0.8897$, the
largest age disparity present; and a subsampling study over the full mortgage
population shows that at the audit size we originally used, a no-floor audit
falsely clears a genuinely disadvantaged subgroup in $83.8\%$ of draws while
a reporting floor never does, because it never reports the subgroup at all
and misidentifies the worst-treated group in $100\%$ of draws. In the audit's
own instruments: a refusal detector that operationalises ``refusal'' as
\emph{under 40 characters} flagged six correct answers and contributed to
permanently blocking a provider, and most of a cross-model ``prompt
sensitivity'' spread came from our own scorer penalising instructed brevity.

The implication is that a fairness verdict is not a property of the model
alone; it is a property of the model plus an unrecorded configuration. The
same data honestly support ``worst-treated group at $0.67$,'' ``at $0.90$,''
or ``unknowable,'' depending on choices no reader can see. An audit
should therefore ship its configuration the way results ship error bars. We
close with the six-item record that would accomplish this, and with
appendices that print the full cell tables and a claim-by-claim map from
every number to the artifact that regenerates it.
\end{abstract}

\begin{CCSXML}
<ccs2012>
<concept>
<concept_id>10010147.10010257</concept_id>
<concept_desc>Computing methodologies~Machine learning</concept_desc>
<concept_significance>500</concept_significance>
</concept>
<concept>
<concept_id>10003456.10003462.10003480</concept_id>
<concept_desc>Social and professional topics~Governmental regulations</concept_desc>
<concept_significance>300</concept_significance>
</concept>
<concept>
<concept_id>10003120.10003121</concept_id>
<concept_desc>Human-centered computing~Human computer interaction (HCI)</concept_desc>
<concept_significance>300</concept_significance>
</concept>
</ccs2012>
\end{CCSXML}

\ccsdesc[500]{Computing methodologies~Machine learning}
\ccsdesc[300]{Social and professional topics~Governmental regulations}
\ccsdesc[300]{Human-centered computing~Human computer interaction (HCI)}

\keywords{measurement, construct validity, algorithmic auditing, disparate
impact, intersectionality, lending, large language models}

\maketitle

\section{Introduction}
\label{sec:intro}

A fairness verdict is not a property of a model alone. It is a property of
the model plus a configuration: which outcome counts as favorable, who
counts as protected, whom a disparity is measured against, how much evidence
the number rests on, and which instruments check the audit's own components.
None of those settings appear in the audit's output, and each is an instance
of one underlying event: the construct a fairness instrument is \emph{named
for} differs from the construct it actually \emph{operationalises}
\citep{jacobswallach2021}, the output is clean either way, and nothing
records the substitution. This paper measures how often each such setting
changes the verdict, and one measured regularity makes the gap dangerous
rather than merely untidy: every wrong verdict erred in the same direction,
toward the clean, reassuring reading.

Lending is the right place to make that measurement. Credit decisions are
among the most consequential automated decisions people
encounter, and two of them dominate household outcomes: consumer installment
credit, which governs access to short-term liquidity, and residential mortgage
origination, which governs access to the principal asset most households will
ever own. \citet{lee2020} make the organising commitment of human-centered
fair AI that fairness is defined relative to a human decision context rather
than in metric space. Lending's decision context is unusually well
documented: the mortgage data are published for public audit, the protected
categories come with written definitions, and the datasets themselves ship the
flags those definitions turn on. That is what turns the gap between named and
operationalised construct from something arguable into something measurable.

The gap appears at every layer of an audit, takes the same shape at each
layer, and is invisible in the output.
The layers are: the metric library, whose defaults resolve normative questions
no caller answered; the mapping from attribute to construct, where the
quantity computed and the quantity the audit's readers care about part
company; and the audit's own evaluation instruments (detectors, gates,
scores), each of which operationalises a judgement (``what counts as an
answer,'' ``what counts as evidence,'' ``what counts as a remedy'') as a
mechanical proxy. At each layer the instrument computes correctly, the output
is clean, and the substitution is recorded nowhere.
Table~\ref{tab:inventory} is the paper in one place: every instance, its
named construct, its operationalisation, and the clean number the gap
produced.

That the instrument is part of the finding is by now well documented.
\citet{deng2022} find practitioners using fairness toolkits in ways their
designers did not anticipate; \citet{leesingh2021} map systematic gaps in what
open-source toolkits support; \citet{holstein2019} report that practitioner
needs diverge sharply from what the literature supplies. \citet{mei2025} show,
in a speech setting, that audit \emph{pipelines} introduce pitfalls that change
conclusions, and \citet{zaccour2025} show that quantitative audits are more
often blocked by data access than by method. We add a setting where the named
constructs come with written definitions and published flags, so the gap can
be measured rather than argued.

\begin{table*}[t]
\caption{The instance inventory. Each row is one measurement in this audit
where the construct operationalised is not the construct named. Every number
is regenerable; appendix~\ref{app:map} maps each to its artifact and command.}
\label{tab:inventory}
\small
\begin{tabular}{@{}p{2.5cm}p{3.6cm}p{4.6cm}p{4.2cm}l@{}}
\toprule
\textbf{Construct named} & \textbf{Operationalised as} & \textbf{The gap} &
\textbf{The clean output} & \textbf{Where}\\
\midrule
\multicolumn{5}{@{}l}{\textit{Layer one: the metric library (verified in AIF360 0.6.1 / Fairlearn 0.13.0)}}\\
\addlinespace[2pt]
``favorable outcome'' & \texttt{favorable\_label=1.0}, a constructor default
& on a \texttt{loan\_default} target, defaulting becomes the good outcome
& $\mathrm{DI}=1.2072$ (disadvantaged group reported as favoured); oriented $0.9862$
& \S\ref{sec:library}\\
\addlinespace[2pt]
``disparate impact'' & a bare point ratio & no interval, $n$, or reference
group ships with the number & $0.6734$ at $n{=}31$ prints as authoritatively
as $0.9604$ at $n{=}4{,}019$ & \S\ref{sec:library}, \S\ref{sec:intervals}\\
\addlinespace[2pt]
``compared to whom'' & highest-rate cell vs.\ designated control, unrecorded
& the convention moves every ratio 3--4 points & $0.6734$ or $0.7012$; $0.8632$
or $0.8988$ & \S\ref{sec:reference}\\
\midrule
\multicolumn{5}{@{}l}{\textit{Layer two: the documented construct (definitions the data and their publishers supply)}}\\
\addlinespace[2pt]
``age'' (protected class) & banded \texttt{age\_group} audited at $40{+}$
& public age bands straddle 62, the boundary the data's own flag encodes;
young and senior pool and cancel & $0.9791$, severity \textsc{low} ($62{+}$
cut: $0.8897$) & \S\ref{sec:age}\\
\addlinespace[2pt]
``discrimination'' & outcome-rate disparity & silent about the model's inputs
entirely & clean $\mathrm{DI}=0.8888$ from a model taking sex and age as
inputs & \S\ref{sec:treatment}\\
\midrule
\multicolumn{5}{@{}l}{\textit{Layer three: the audit's own instruments (our released harness; evidence about building an audit)}}\\
\addlinespace[2pt]
``refusal'' & any narrative field under 40 characters & instruction-following
terseness is flagged; API failures conflated with model behaviour & 6 of 9
flagged with zero refusals; a provider blocked at a measured $19.79$
& \S\ref{sec:refusal}\\
\addlinespace[2pt]
``fabrication'' & a float literal absent from context & derivable arithmetic
($0.8888 \to 88.88\%$) counts as fabrication & the model that shows its work
flags $17$ vs.\ $4$ & \S\ref{sec:fabrication}\\
\addlinespace[2pt]
``mitigation'' & a canonical algorithm name & where the problem is data
scarcity, ``collect more data'' is the only defensible answer & the gate fails
exactly those two attributes & \S\ref{sec:mitigation}\\
\addlinespace[2pt]
``severity'' & \textsc{critical} below $0.72 = 0.80 \times 0.9$ & the boundary
has no statutory or statistical provenance & authoritative-looking band labels
& \S\ref{sec:ownthresholds}\\
\addlinespace[2pt]
``capability'' & mean score over heterogeneous tasks & the mean tracks task
difficulty; the scorer penalises terse output & $75\%$ agreement decomposing
to $33.3/91.7/100$; a $14.7$-point ``prompt sensitivity'' spread that is $6.7$
without the terse cycle & \S\ref{sec:capability}\\
\bottomrule
\end{tabular}
\end{table*}

\paragraph{Contributions.}
(i)~Three library-layer gaps verified directly in AIF360~0.6 and
Fairlearn~0.13 and scoped to what we verified: a favorable-label constructor
default that inverts a ratio's meaning, the absence of any uncertainty
quantification on disaggregated ratios, and the absence of any record of the
reference group (\S\ref{sec:library}, appendix~\ref{app:repro}).
(ii)~Two definition-layer gaps measured on real lending data: the banded-age
operationalisation cannot represent the $62{+}$ boundary the data's own flag
encodes, and outcome-rate tooling is structurally silent about the model's
inputs (\S\ref{sec:statute}).
(iii)~A subsampling study ($2{,}000$ draws at each of five audit sizes,
five reporting policies, three reference conventions) that converts this
audit's founding anecdote into a rate: policies that report small cells
falsely clear violating subgroups in $84$--$95\%$ of draws at the audit size
we originally used, and policies that suppress them never name the
worst-treated group at all (\S\ref{sec:subsample}, appendix~\ref{app:subsample}).
(iv)~Five instrument-layer gaps in our own released harness, reported as
evidence about building an audit, including a reanalysis showing that a
cross-model prompt-sensitivity difference we previously reported is partly an
artifact of how our scorer treats terse output (\S\ref{sec:instruments}).
(v)~Full-transparency appendices: every cell table under every convention,
the minimal reproduction, and a claim-by-claim artifact map enforced by a
test suite (appendices~\ref{app:cells}--\ref{app:map}).

\section{Related Work}
\label{sec:background}

\paragraph{Measurement.} \citet{jacobswallach2021} supply the vocabulary this
paper instantiates: fairness quantities are \emph{measurement models} of
unobservable theoretical constructs, and the question to ask of any number is
whether the operationalisation matches the construct it claims to measure,
which is construct validity. \citet{passi2019} locate the same gap upstream,
in problem
formulation; \citet{selbst2019} describe the failure of porting formal
criteria across social contexts; \citet{obermeyer2019} give the canonical
in-vivo case, a proxy (cost) silently standing in for a construct (need).
Disaggregated evaluation has its own choices (which groups, which cells,
which comparisons) that \citet{barocas2021} catalogue; \citet{chen2018}
supply the bias--variance decomposition that governs what a subgroup estimate
can support, the statistical half of our \S\ref{sec:validity}. On
intersectionality, \citet{crenshaw1989} names the construct our
\texttt{race}$\times$\texttt{sex} cells operationalise; \citet{buolamwini2018}
give the canonical empirical audit, \citet{foulds2020} a formal definition,
and \citet{kearns2018} the auditing-theoretic treatment. That fairness
criteria conflict is settled \citep{chouldechova2017,kleinberg2017,hardt2016};
\citet{corbettdavies2018} catalogue how measures mislead when detached from
decision context.

\paragraph{The documented decision context.} We make no legal claims in this
paper: discrimination law does not reduce to a metric threshold
\citep{barocas2016,wachter2021}, and adjudicating it is not our competence.
What lending supplies, and what this paper uses, is narrower: the decision
context comes with \emph{written definitions and published flags}. The public
mortgage data ship a dedicated \texttt{applicant\_age\_above\_62} field
because $62{+}$ is a defined boundary in that domain's own documentation, so
whether an age operationalisation can represent that boundary is a checkable
question of measurement, not of law. Likewise, whether a model consumes a
protected attribute directly is a documented property of its inputs
(\citet{gillis2022} examines input scrutiny at length), and no outcome-rate
metric reports it (\S\ref{sec:treatment}).

\paragraph{The audit's instruments.} The third layer is newer but not new.
LLM-as-judge evaluation is known to be fragile \citep{zheng2023};
\citet{wester2024} study how models deny requests, the construct our
refusal detector was built to measure; \citet{sclar2024} show scores move
under semantically null formatting changes, the construct our
prompt-sensitivity comparison was built to measure. What this literature and
the measurement literature have not yet been made to say jointly is that a
detector, a gate, or a score \emph{is} a measurement model, with the same
validity obligations as the fairness metric it polices. Documentation and
process interventions \citep{mitchell2019,gebru2021,madaio2020} share our
premise that fairness work needs infrastructure; audit practice
\citep{metaxa2021,raji2019,raji2020} and traceability \citep{kroll2021} give
the released-artifact design its rationale. The stakeholder line
\citep{luo2025,luo2026,luo2026b,deshpande2022,stumpf2024,kim2025} asks who
should choose the criteria; our contribution is adjacent, showing that the
apparatus resolves several such questions before anyone is asked, at layers
where nobody thinks to look. \citet{zamfirescu2022} diagnose human error
inside an analysis pipeline; we diagnose the pipeline's own rulers.

\section{Methods}
\label{sec:apparatus}

\subsection{Settings and data}
Two regulated lending decisions and one instrument-test
setting. German Credit \citep{germancredit} stands in for consumer installment
credit: $1{,}000$ applications of 1990s German data, a structural analogue
from a different market and era, from which we draw no US conclusion. HMDA
Georgia
\citep{hmda} is residential mortgage origination: $10{,}978$ applications
after population construction (below). The third setting, a Lending Club
default task, contains \emph{no protected class}: its \texttt{gender} column
is synthetic, assigned by an independent coin flip
(\texttt{np.random.random(n) > 0.5}). We retain it because a harness must be
tested somewhere its verdicts are known to be vacuous, and two of this paper's
instrument findings (\S\ref{sec:instruments}) surfaced there; nothing in this
paper treats it as a lending sector. Lending discrimination itself is
documented across this span \citep{duarte2012,pope2011,bartlett2019,butler2017,sarwal2024}.

\subsection{Deterministic metric layer}
A pipeline computes disparate impact with
AIF360 \citep{aif360}, cross-checked against Fairlearn \citep{fairlearn}, and
assigns severity labels by a fixed rule. The per-attribute summary the
downstream layers receive carries three fields (\texttt{DisparateImpact},
\texttt{TheilIndex}, \texttt{BiasFlag}), of which the Theil index is
computed on the prediction vector alone and is therefore identical across
attributes within a dataset; the per-dataset CSV carries difference metrics
(demographic parity, equal opportunity, average odds) whose full-population
values appear in appendix~\ref{app:map}'s artifacts.\footnote{The Lending Club
summary uses a different five-field schema than the two regulated settings,
a dual-schema defect we report rather than repair, since repairing it would
churn every downstream artifact.} Population, not sample: every applicant is
scored out-of-fold by stratified five-fold cross-validation (random forest,
accuracy $0.8002$, AUC $0.8018$ on HMDA), so audits run over all $10{,}978$
and $1{,}000$ rows rather than a held-out split.

\subsection{LLM interpretation layer}
A language model interprets the computed
metrics and never recomputes one; scoring is against the deterministic label,
never another model's opinion, with one disclosed exception
(\S\ref{sec:capability}). Track~Q issues four frozen prompt strategies per
(setting, attribute) pair and scores the returned severity against the
deterministic label; Track~H produces narrative audits judged against
mechanical checks. Prompt builders, detectors, severity rule and rubric are
printed verbatim in appendix~\ref{app:instruments}; a no-key replay path
reproduces all scoring.

\subsection{Population construction}
The HMDA file begins at $20{,}000$ rows
and the audit runs on $10{,}978$: an attrition of $45.1\%$ that includes
dropping \texttt{Race Not Available} ($4{,}140$ rows, $20.7\%$ of the file),
joint applications, and free-form race text. That is selection on the
protected attribute, performed before any metric runs; \S\ref{sec:invisible}
discusses it and appendix~\ref{app:attrition} prints the rule-by-rule table
with the disparate-impact ratios recomputed with the dropped categories
retained. HMDA's public data carry no credit score, the primary underwriting
variable.

\subsection{The four-fifths screen}
One convention in this paper is itself an imported construct, and we label it
rather than launder it. The four-fifths screen originated in US employment
guidance, not in credit; we retain $0.80$ because it is the de facto screening
convention in fairness practice, we use it only as a screen, and we note that
the guidance's own text contains the caveat our \S\ref{sec:validity}
quantifies: greater differences ``may not constitute adverse impact where the
differences are based on small numbers and are not statistically
significant.'' The screen's own source anticipates the failure mode this
paper measures.

\section{Results}
\label{sec:results}

Results are ordered by layer, following Table~\ref{tab:inventory}: the metric
library (\S\ref{sec:library}), the mapping from attribute to documented
construct (\S\ref{sec:statute}), the statistical validity of what either
layer reports (\S\ref{sec:validity}), and the audit's own instruments
(\S\ref{sec:instruments}).

\subsection{Layer one: library defaults}
\label{sec:library}

\emph{Construct named: ``favorable outcome,'' ``disparate impact.''
Operationalised as: a constructor default of $1.0$; a bare point ratio.}

We inspected AIF360~0.6.1 and Fairlearn~0.13.0 directly; the claims in this
section are about those libraries as verified, not about the field.

\paragraph{Which outcome is good is a default.} AIF360's
\texttt{BinaryLabelDataset} takes \texttt{favorable\_label} with a default of
$1.0$. Credit risk targets are conventionally coded with the adverse event as
$1$ (default, delinquency, charge-off, fraud), so on a column named
\texttt{loan\_default} the default asserts that defaulting is the good
outcome. The minimal reproduction (appendix~\ref{app:repro}; deterministic,
pinned versions) returns $\mathrm{DI}=1.2072$ under the library default, with
the disadvantaged group reported as \emph{favoured}, and $0.9862$ under
correct orientation. The scope of this claim matters and we state it
precisely: this is one constructor in a two-API library. AIF360's other entry
point, \texttt{StandardDataset}, makes \texttt{favorable\_classes} a required
argument and cannot be constructed without answering the question. Fairlearn's
\texttt{demographic\_parity\_ratio} exposes no favorable-label parameter at
all; it computes a min/max ratio of the prediction column \emph{as encoded},
so on the same table it reports $0.8284$, a ratio of default rates and a
number about the wrong outcome entirely, and delegates the orientation
question to whoever encoded the column, without recording that it did so.

The in-vivo version ran end to end in our own harness before we caught it:
selection rates computed on predicted default gave the Lending Club synthetic
\texttt{gender} attribute $\mathrm{DI}=0.0$ and a \textsc{critical} label,
the most severe output the pipeline can produce, where the oriented value
is $0.9931$, near parity. Nothing was misused; AIF360 computed what it was
asked. A target named \texttt{approved} makes label $1$ favorable, a target
named \texttt{loan\_default} inverts it, and nothing in the metric signature
records which case holds. It is now a required argument and a regression test
in our pipeline, which converts a silent inversion into a build failure.

\paragraph{No uncertainty ships with a disaggregated ratio.} Neither library
attaches an interval or a standard error to a disparate-impact ratio, and
neither reports the $n$ it was computed on unless asked (Fairlearn offers
\texttt{count} as an opt-in metric). A ratio resting on $31$ people prints
with the same four decimal places as one resting on $4{,}019$. This is the
defensible core of what practitioners experience as the ``minimum subgroup
size'' problem: a floor is arguably not a metric library's job, but
uncertainty is what the estimate \emph{means}, and \S\ref{sec:validity}
measures what its absence costs.

\paragraph{Nothing records the comparison.} A disparate-impact ratio has a
denominator, and neither library's output records which group supplied it.
\S\ref{sec:reference} shows the choice moves every number in our tables.

\subsection{Layer two: the documented construct}
\label{sec:statute}

\subsubsection{The protected-class boundary}
\label{sec:age}

\emph{Construct named: age, as the lending domain defines the protected
boundary. Operationalised as: a banded \texttt{age\_group} attribute audited
at $40{+}$.}

Fairness work on credit routinely audits age at $40{+}$, and our pipeline did,
on the strength of a code comment asserting that the threshold aligned with
the domain's definition. It does not: $40{+}$ is a threshold imported from
employment practice. The boundary lending's own documentation defines is
$62{+}$, which is why the public mortgage data ship a dedicated
\texttt{applicant\_age\_above\_62} flag in the first place.

The substitution is not cosmetic. Our banded age attribute pools young and
senior applicants ($n{=}2{,}304$ and $3{,}167$) against the middle band; the
two move in opposite directions, so pooling cancels them. The banded attribute
returns $\mathrm{DI}=0.9791$, severity \textsc{low}: a clean bill of
health. Audited on the documented cut, the same model over the same $10{,}978$
rows returns $0.8897$, the largest age disparity in the dataset and the lowest
of its four protected attributes. The flag required was present and unused in
the raw HMDA file, and it is necessary rather than convenient: the public age
bands straddle $62$ (the ``55--64'' band cannot be split), so \emph{no} banded
attribute can represent the boundary. No re-binning fixes a partition whose
cuts are in the wrong places. We report honestly that $0.8897$ clears the
$0.80$ screen and the pipeline's \texttt{BiasFlag} remains false on both cuts:
the boundary changes which question the number answers, not the verdict on
this data.

\subsubsection{The comparator}
\label{sec:reference}

\emph{Construct named: disparity against the favoured. Operationalised as: a
ratio against an unrecorded, convention-chosen denominator.}

Every ratio in this paper's tables is measured, by our pipeline's convention,
against the most-favoured adequately-sized cell, and no metric signature
records that. Measured instead against a designated control group
(\texttt{White $\times$ Male}, the comparison conventional in US fair lending
analysis), the same cells read differently: \texttt{American Indian $\times$
Male} $0.7012$ rather than $0.6734$, \texttt{Black $\times$ Female} $0.8988$
rather than $0.8632$. The convention moves every number in the table by three
to four points (appendix~\ref{app:cells} prints all cells under three
conventions), and \S\ref{sec:subsample} shows the choice also changes how
often an audit is \emph{wrong}, because a data-dependent reference is itself
re-estimated in every sample.

\subsubsection{The inputs the metric cannot see}
\label{sec:treatment}

\emph{Construct named: discrimination. Operationalised as: outcome-rate
disparity.}

The German Credit model takes applicant sex and age as input features
(direct use of the protected attributes, the first thing any fairness review
of a lending model would ask about), and our pipeline ran that
specification end to end and emitted a clean $\mathrm{DI}=0.8888$ without
comment. We previously filed this under limitations; it is a result, and
arguably the paper's sharpest. An outcome-rate metric is a function of
predictions and group labels alone: it is structurally silent about what the
model consumes, so the cleanest possible metric output is compatible with the
most direct possible use of a protected attribute. No disparity ratio, at any
sample size, under any convention, could have flagged it. Input scrutiny
and outcome scrutiny are different measurements \citep{gillis2022}, and a
toolchain that computes only the second says nothing about the first.

\subsection{Measurement validity}
\label{sec:validity}

\subsubsection{Intervals on the real audit}
\label{sec:intervals}

Our audit originally ran on a held-out split of $2{,}196$ mortgage
applications and reported, among other things, that American Indian men were
the \emph{best}-treated group in the data ($\mathrm{DI}=1.1504$), on the
evidence of four applicants, who all happened to be approved. Over the full
population the same group is the worst cell in the dataset: $0.6734$ at
$n{=}31$. The audit did not miss the disparity; it \emph{inverted} it, and
attached a reassuring number to the inversion. The floor was right about what
it suppressed: \texttt{American Indian $\times$ Female} read $0.1917$ on the
split (one approval in six), which looks like the most severe disparity in the
study and is noise: $0.8815$ at $n{=}25$ in full. Small-group estimates are
unreliable in both directions, and every affected cell belonged to a racial
minority, which is mechanism rather than coincidence: a floor on group size is
a floor on group population, so the groups whose treatment is least certain
are exactly the groups least likely to be counted.

Table~\ref{tab:cells} prints the full-population
\texttt{race}$\times$\texttt{sex} table, in full because this paper's
argument is precisely that partial tables mislead, with the Wilson-interval
disparate-impact bounds beside every ratio (reference: \texttt{Asian $\times$
Female}, $n{=}254$, the pipeline's most-favoured-adequate-cell convention;
floor $15$). Six of twelve cells have no determinate severity band: their
intervals straddle a band boundary, so the audit does not know which band they
belong to. \texttt{American Indian $\times$ Male}, the lowest point estimate,
has interval $[0.453, 0.904]$: consistent with a critical violation and
with clearing the screen. It cannot be called the worst-treated group. What
survives is narrower and firmer: \texttt{Black $\times$ Female} ($0.8632$,
$[0.806, 0.937]$, $n{=}1{,}994$) and \texttt{Black $\times$ Male} ($0.8643$,
$[0.805, 0.940]$, $n{=}1{,}707$) have intervals wholly below parity. The
defensible reading of this dataset is a real, precisely estimated disparity
against Black applicants of roughly $0.86$, which a point-estimate ranking
demotes to sixth and seventh place, behind five cells whose ordering is noise.
Reporting a ratio to four decimals without an interval is what makes that
demotion possible \citep{chen2018}.

\begin{table}[t]
\caption{HMDA \texttt{race}$\times$\texttt{sex}, full population
($10{,}978$), out-of-fold predictions. Reference \texttt{Asian $\times$
Female} ($n{=}254$); conservative Wilson-based DI intervals; ``band det.''
marks cells whose whole interval sits in one severity band. Generated by
\texttt{scripts/interval\_estimates.py}; appendix~\ref{app:cells} repeats this
table under two further reference conventions.}
\label{tab:cells}
\footnotesize
\input{generated/appendix_a_hmda_max_rate_floor}
\end{table}

\subsubsection{The subsampling study}
\label{sec:subsample}

The anecdote above is one draw from a distribution, so we measured the
distribution. Treating the full-population table as ground truth $\theta$, we
drew $B{=}2{,}000$ simple random subsamples at each audit size $m \in \{500,
1{,}000, 2{,}196, 4{,}000, 8{,}000\}$ and ran five reporting policies on every
draw: a floor at $n \geq 15$ (the pipeline's), a floor at $n \geq 30$ (which
our own codebase also contained, in an adjacent module), no-floor point
estimates, no-floor Wilson intervals that assert a verdict only when the
interval clears or fails the screen entirely, and empirical-Bayes shrinkage
toward the pooled rate. Each policy runs under three reference conventions,
with $\theta$ defined per convention. We record how often each policy asserts a
population-violating cell clear (false clearance), asserts a clearing cell in
violation (false alarm), places a reported estimate on the wrong side of
parity (sign inversion), and names the wrong worst-treated cell or declines to
name one. Appendix~\ref{app:subsample} gives the full grid, the estimator
specification, and the design notes; Figure~\ref{fig:subsample} and
Table~\ref{tab:subsample} give the primary convention.

\begin{figure*}[t]
\centering
\includegraphics[width=\textwidth]{generated/subsample_rates.pdf}
\caption{Error rates by audit size $m$ ($B{=}2{,}000$ draws per size, shaded
bands are $95\%$ Wilson intervals on the rate; primary reference convention).
Left: how often each policy asserts a population-violating cell clear of the
$0.80$ screen. The floor and interval policies coincide at zero for
$m \leq 2{,}196$ (floors by silence, intervals by abstention), and the
$n\geq15$ floor begins to leak exactly when the violating cells cross it
($26\%$ at $m{=}8{,}000$). Right: how often the policy names the wrong
worst-treated cell, given that it names one.}
\label{fig:subsample}
\end{figure*}

\begin{table*}[t]
\caption{All five policies at $m{=}2{,}196$, the size of the audit we
originally ran. Rates are \% of draws with $95\%$ Wilson intervals;
``abstains'' is the share of draws where the policy declines to name a worst
cell; ``cells suppressed'' is the mean count of the twelve
\texttt{race}$\times$\texttt{sex} cells not reported at all.}
\label{tab:subsample}
\small
\input{generated/subsample_body_table}
\end{table*}

Three results order the section. First, at the audit size we actually used,
\emph{every policy that reports small cells is usually wrong}: no-floor point
estimates falsely clear a violating cell in $83.8\%$ of draws and put at least
one reported estimate on the wrong side of parity in $79.0\%$. Our
$1.1504$ was not bad luck. Second, the floors never false-clear at this size
for the least reassuring possible reason: they never report the violating
cells at all. Both floors suppress six of twelve cells on average and
misidentify the worst-treated group in $100\%$ of draws; and the $n \geq 15$
floor starts leaking at $m{=}8{,}000$ ($26.4\%$ false clearance), exactly when
the violating cells begin to cross it with noisy estimates. A floor does not
solve the small-cell problem; it postpones it to a larger audit and converts
it to silence in the meantime. Third, the two policies that respect
uncertainty fail in opposite, honest ways: the Wilson policy never
false-clears and almost never misranks ($6.0\%$), but abstains from naming a
worst cell in $90.3\%$ of draws (at full population its interval on the
worst cell still straddles the screen), while empirical-Bayes shrinkage is
the \emph{most} confident false-clearer of all ($95.2\%$ at $m{=}2{,}196$,
rising with $m$), because a prior centered on the pooled rate absorbs exactly
the small cells where the disparities live. Floors trade false clearance for
silence; intervals trade decisiveness for honesty; shrinkage trades variance
for a bias that points, in this data, toward clearance. The choice among them
is a policy about who becomes visible, and nothing in a reported ratio records
which policy produced it.

\subsection{Layer three: audit instruments}
\label{sec:instruments}

The findings in this section are about our released harness, measured on
frozen artifacts, and are evidence about \emph{building} an audit rather than
about the field. They are in this paper because they are the same finding as
\S\ref{sec:library} and \S\ref{sec:statute}: an instrument named for one
construct, computing another, cleanly. Appendix~\ref{app:instruments} prints
every implementation verbatim.

\subsubsection{``Refusal'' := fewer than 40 characters}
\label{sec:refusal}

\emph{Construct named: refusal to answer \citep{wester2024}. Operationalised
as: any narrative field under 40 characters.}

The detector flagged six of nine responses under the terse-output prompt
strategy. None was a refusal: no flagged record contained a refusal phrase, a
missing field, or an invalid severity. The detector treats any narrative field
under 40 characters as effectively empty, so \texttt{how\_to\_fix:
"DisparateImpactRemover"} (correct, on-spec, and naming the applicable
post-processor \citep{feldman2015} in 22 characters) registered as a
refusal. Compliance with the instruction was the trigger. Reporting the raw
output would have yielded a $66.7\%$ refusal rate and a tidy story about
constrained prompting inducing evasion, entirely an artifact of our
instrument. (Three of the six flags are Lending Club records: the harness
failed identically on the setting whose verdicts are vacuous, which is what
that setting is for.)

Then it set policy. Applied to a second provider, the same detector flagged 19
of 24 cycles. Reanalysis of the frozen records
(appendix~\ref{app:instruments}) separates what the flag conflated: 18 of the
24 records are \emph{empty API results} (transport failures, not model
behaviour), of which 14 scored zero and four scored $10.0$ despite containing
nothing, because the cause-alignment subscore returns a flat $10.0$ whenever
the deterministic baseline lists no cause labels. The provider's mean is
$19.79$ over all records; $47.50$ excluding zero scores (the figure we
previously reported, which still contains the four empty-but-scored records);
and $72.50$ over the six records that contain any content at all. One of
those six, a terse correct answer scoring $85$, is the only content-bearing
record the refusal detector flagged. On the strength of the $19.79$ figure, that
provider was recorded as permanently blocked in our guardrail configuration,
where the block persists. An invalid measurement did not merely mislead a
table; it produced a durable governance decision about a vendor.

\subsubsection{``Fabrication'' := a float literal absent from context}
\label{sec:fabrication}

\emph{Construct named: asserting numbers that were never provided.
Operationalised as: any decimal literal not string-matching the context.}

Replaying identical packs against two model generations,
\texttt{gpt-5-mini} trips $17$ flags across $28$ calls where \texttt{gpt-4o}
trips $4$. Read naively, the newer model fabricates five times as often.
Checking each flagged value against the context, $14$ of $21$ ($67\%$) are
derivable from it: $88.88$ is the disparate impact $0.8888$ restated as a
percentage, $83.87$ is $0.8387$, and most of the remainder are gaps between
two provided values expressed in percentage points. The detector matches raw
float literals, so converting a ratio to a percentage or subtracting two given
numbers registers as fabrication. The model that reasons more explicitly is
ranked as the less trustworthy one, precisely because it shows its work. The
seven values that are not derivable ($19.9$ three times, $5.3$ four times)
remain honest fabrication candidates, and we report them as such.

\subsubsection{``Mitigation'' := a canonical algorithm name}
\label{sec:mitigation}

\emph{Construct named: a specific, actionable remedy. Operationalised as: a
keyword match against a list of algorithm names.}

A guardrail requires every proposed mitigation to name a canonical algorithm:
reweighing \citep{kamiran2011}, disparate impact removal
\citep{feldman2015}, and similar. It failed exactly the two attributes whose
problem \emph{is} data scarcity, where the audit proposed acquiring records to
a stated floor, hierarchical shrinkage, and manual review. No post-processor
repairs a six-record subgroup. The gate encodes an assumption that every
fairness problem has an algorithmic mitigation, and where that is false it
penalises the only defensible answer.

\subsubsection{Our own thresholds, same class}
\label{sec:ownthresholds}

\emph{Construct named: severity. Operationalised as: fixed cuts at $0.72$ and
$0.80$.}

The pipeline's severity rule places its \textsc{critical} boundary at $0.72$.
That number is $0.80 \times 0.9$, constructed in this project, with no
external provenance, regulatory or statistical; it looks authoritative for
the same reason the mis-sourced $40{+}$ code comment did. The same codebase
carried two reporting floors, $n \geq 15$ and $n \geq 30$, in adjacent
modules, so the same cell was reportable or not depending on which module
read it; \S\ref{sec:subsample} shows those two floors behave identically at
small $m$
and diverge exactly at the audit sizes where the $n \geq 15$ floor starts to
leak. This is why the analysis in \S\ref{sec:validity} uses only the labeled
$0.80$ screen plus intervals, and quotes the pipeline's band labels as
pipeline convention rather than as findings. The code is left as implemented;
the exhibit is the point.

\subsubsection{Scores that operationalise ``capability''}
\label{sec:capability}

\emph{Construct named: how good the model layer is. Operationalised as: mean
score over heterogeneous tasks, graded by mechanical subscores.}

\paragraph{An agreement mean tracks the difficulty mix.} Severity agreement
between \texttt{gpt-4o} and the deterministic labels is $27/36 = 75.0\%$. It
decomposes to $33.3\%$, $91.7\%$ and $100.0\%$ across the three settings, and
the ordering tracks how discriminating each audited classifier is rather than
how capable the model is: the $100\%$ case is the setting where the classifier
barely separates anyone, so every attribute is a trivial near-parity
\textsc{low} and agreeing costs nothing. A single headline agreement figure
for an LLM auditor is not a well-defined quantity; in general, any benchmark
averaging over heterogeneous instances reports a number dominated by its
difficulty mix. Errors were asymmetric in the safe
direction (seven over-escalations, two under-escalations), which matters
because under-escalation is the consequential failure for an audit tool.

\paragraph{Track H, disclosed.} The narrative-audit track's rubric scores are
one model's opinion of another's work (\texttt{gpt-4o} grading narratives
generated by \texttt{o3}), which is exactly what this paper's architecture
forbids: scoring against the deterministic label, never another model's
opinion. We exclude those scores ($4.8$, $4.7$, $5.0$ of $5$) rather than
report them; an LLM-judged ``quality'' score is itself an unvalidated
operationalisation \citep{zheng2023}, and this one was ours. What survives on
mechanical evidence: all $27$ of $27$ citations in the narrative audits
verified against the retrieved evidence pool. Two behaviours worth reporting
as observations rather than scores: on a six-record subgroup the narrative
audit independently reproduced the small-$n$ caveat and quantified it (one
label change would swing the ratio below the screen), and on the
near-degenerate classifier it declined the flattering reading.

\paragraph{Prompt sensitivity, reanalysed.} A widely reported result about LLM
evaluators is that prompt phrasing moves the verdict \citep{sclar2024}. It
reproduced here with a generational asymmetry: across four prompt strategies
on identical packs, the spread in cycle-mean score is $14.7$ and $9.0$ points
for \texttt{gpt-4o} (German Credit, HMDA) against $0.0$ and $1.9$ for
\texttt{gpt-5-mini}. We previously reported that contrast as
model-generational. Reanalysis of the frozen records
(appendix~\ref{app:instruments}) shows the refusal flag never feeds the score,
so no refusal penalty inflates it. But the \texttt{gpt-4o} spread is
concentrated in the one strategy that demands terse output: excluding the
constrained cycle it falls to $6.7$ and $1.25$, and the cycle-four dip
decomposes almost entirely into the completeness and severity-agreement
subscores ($78\%/22\%$ on German Credit, $42\%/58\%$ on HMDA), subscores
that reward enumerating content, scored on answers instructed to be short.
Every \texttt{gpt-4o} constrained-cycle record was simultaneously flagged by
the 40-character refusal rule; the scorer and the detector were measuring the
same thing, brevity, under two different construct names. A generational
difference survives (\texttt{gpt-5-mini}'s spread is near zero under every
treatment), but the published magnitude was partly our instrument, and
paraphrase-instability claims should carry a model, a date, \emph{and} an
instrument audit. (\texttt{gpt-5-mini} also rejects the \texttt{temperature}
parameter outright, so a repeated-sampling study is impossible on it, an
API constraint that silently determines which experiments a benchmark can
run.)

\section{Discussion}
\label{sec:invisible}

The layer-one findings hold for anyone using the verified libraries; the rest
are about choices our audit made because nothing made them for us, and that
is the point. An audit's numbers mean what its readers take them to mean only
if the construct questions are answered correctly: which outcome is favorable,
who counts as protected, against whom, on how much evidence, checked by which
instruments.
The toolchain answers the first by default and the rest not at all, so each is
settled privately, by whoever writes the harness, without a record.

The asymmetry that makes this dangerous is directional. A false alarm costs
reviewer time; a false clearance closes an open question. Nearly every gap in
Table~\ref{tab:inventory} produced its error in the clearance direction: the
banded age attribute read \textsc{low}, the small-sample audit read its worst
cell as its best, the shrinkage policy cleared violating cells with the most
confidence of any policy, and the ``prompt sensitivity'' figure flattered the
older model's critics rather than its scorer. Clean results are the outputs
least likely to be re-examined, and every mechanism in this paper pushes
toward clean.

The instruments make it worse, because every mechanical proxy for judgement is
a policy. A length threshold that treats brevity as refusal is a policy about
what counts as an answer. A keyword list that treats ``acquire more data'' as
a non-mitigation is a policy about what counts as a remedy. A float-literal
match that treats a percentage as a fabrication is a policy about what counts
as evidence. None was labelled a policy; each was a heuristic written to make
an evaluation automatable, and each punished the better answer in the cases
where judgement was most needed. As evaluation automates further (LLM
judges, automated red-teams, compliance scoring), this class of failure
scales with it.

Two of our own decisions belong in this catalogue, and we close the section
with them rather than with other people's tools. Our population construction
dropped $45.1\%$ of the raw HMDA file before any metric ran, including every
applicant whose race is not available ($20.7\%$ of the file): selection on
the protected attribute, deciding who counts as protected upstream of every
number in this paper. Appendix~\ref{app:attrition} prints what the dropped
categories look like when retained (the excluded \texttt{Race Not Available}
group approves at $0.7234$, below every retained majority group). And our
division of labour (deterministic code owns every number, the model owns
only interpretation, scoring runs against the numbers) protected every
\emph{number} in this paper and not one of the \emph{operationalisations}: it
is the reason we can state which of our results were artifacts, and it is no
defence against choosing the wrong construct in the first place.

\section{Recommendations}
\label{sec:protocol}

Each item follows from a measured failure above; together they are the
construct documentation a reader would need for the numbers to mean what they
appear to mean.

\begin{enumerate}
  \item \textbf{State the favorable outcome as an explicit, required
        argument.} Never infer it from label values. One parameter converts a
        class of silent inversions into a build failure
        (\S\ref{sec:library}).
  \item \textbf{Report the interval, the $n$, and the reference group beside
        every ratio}, and print what any floor withheld: suppressed cells
        with their $n$ and their unreported value. The audit of
        \S\ref{sec:intervals} was not wrong about any number it computed; it
        was wrong about what it did not print.
  \item \textbf{Bind protected-class definitions to their documented source}
        in code, with the citation next to the threshold. A comment asserting
        the wrong provenance survived review precisely because it looked
        authoritative (\S\ref{sec:age}).
  \item \textbf{Treat every mechanical gate as a hypothesis to falsify},
        starting with inputs where the correct answer is unusual: the terse
        answer, the tiny subgroup, the model that shows its work
        (\S\ref{sec:instruments}).
  \item \textbf{Audit the population, not a convenience sample, and disclose
        population construction.} Most of our small-cell instability came from
        auditing a split; all of our attrition happened before any metric ran
        (\S\ref{sec:subsample}, appendix~\ref{app:attrition}).
  \item \textbf{Name the reporting policy.} Floor, interval, or shrinkage is a
        choice among failure modes (silence, abstention, or confident
        clearance), and \S\ref{sec:subsample} shows they disagree massively
        at realistic audit sizes. A reported table should say which policy
        produced it.
\end{enumerate}

\paragraph{These steps are necessary, not sufficient.} The list is what our
measured failures imply, not an enumeration of the ways an audit can be wrong,
and completing it does not make a model fair or a lender compliant. A real
review involves far more than a disparity screen: significance testing,
controls for legitimate underwriting factors, comparative file review, and
analyses of pricing and of less discriminatory alternatives, none of which a
disparity ratio addresses. We perform the screening step only. Anyone citing
this protocol as evidence of compliance is misusing it.

\section{Limitations}
\label{sec:limitations}

\textbf{Scale.} Two lending datasets, one geography and era each, two
primary models. These are mechanism-level findings: we show that a class of
failure occurs and how it operates, not how often it occurs in the wild. The
German data are a structural analogue from a different market and era, and no
US conclusion is drawn from them.

\textbf{No human subjects.} We make claims about what an audit reader would
conclude without testing them on audit readers. That is the natural next
study, and the main gap between this work and the stakeholder line it builds
on \citep{luo2025,luo2026b}.

\textbf{No causal or legal claims.} All metrics are associational; root-cause
mapping is rule-based inference over metric values. Regulation appears in this
paper only as the provenance of definitions and flags, never as something
we interpret, and a disparate-impact ratio below $0.80$ is a screening
signal, not a finding of discrimination \citep{barocas2016,wachter2021}.

\textbf{Model limitations.} The German Credit classifier uses protected
attributes as features (reported as a result, \S\ref{sec:treatment}); the
Lending Club protected attribute is synthetic; HMDA lacks credit score. Each
bounds what the corresponding audit can mean.

\textbf{Sampling confound.} Temperature and seed are unpinned in the
interpretation layer, so within-model cycle differences mix prompt strategy
with sampling noise; the cross-model contrast in \S\ref{sec:capability} is
robust to this only because both models face the same confound, and the
instrument decomposition there is the reason we no longer quote the
within-model spread as a finding. An earlier ungrounded-vs-grounded citation
comparison straddles a pipeline revision and is not quoted anywhere in this
paper because it is a legacy contrast, not an ablation.

\textbf{The subsampling study measures our pipeline's conventions.} Ground
truth is the full-population table under a stated reference convention;
different conventions define different $\theta$, which is part of the finding
but also a limit on transportability. The DI intervals are deliberately
conservative and over-cover ($0.993$--$1.0$ measured against nominal $0.95$);
appendix~\ref{app:subsample} details this and the estimator choices.

\section{Conclusion}

At the layer of the library, a constructor default decided which outcome of a
loan is good, and reported the disadvantaged group as favoured. At the layer
of the definition, an attribute banding decided which age question the audit
asked, and answered the wrong one cleanly; and at realistic audit sizes, the
reporting policy (floor, point, interval, shrinkage) decided whether a
violating subgroup was falsely cleared, silently omitted, or honestly left
undetermined, at rates between $84\%$ and $95\%$ for the policies that report
small cells. At the layer of the audit's own instruments, a 40-character rule
decided what counted as a refusal and blocked a vendor; a float match decided
what counted as fabrication and demoted the model that showed its work; and
two mechanical subscores, rewarding enumeration on answers instructed to be
terse, manufactured most of a ``prompt sensitivity'' result we ourselves
previously reported.

Every one of these is the same event: an intended, documented construct,
operationalised as something adjacent, producing a clean number, recorded
nowhere. The gap is not closable in general, because every measurement
operationalises, but it is \emph{recordable}, and \S\ref{sec:protocol} is
the record this audit would have needed: the outcome, the interval, the $n$,
the reference group, the definition's source, the reporting policy, and an
adversarial test for every gate. None of it is technically difficult. All of
it is currently optional, which is the problem.

\begin{acks}
We thank the maintainers of AIF360 and Fairlearn. Retrieved literature was
obtained via the Semantic Scholar API.
\end{acks}

\bibliographystyle{ACM-Reference-Format}
\bibliography{refs}

\appendix

\section{Full cell tables under three reference conventions}
\label{app:cells}

All tables are generated by \texttt{scripts/interval\_estimates.py --emit-tex}
from the out-of-fold prediction files; nothing is transcribed by hand. The
conventions: \emph{max rate above floor} (the pipeline's rule; reference is
the highest-rate cell with $n \geq 15$), \emph{max rate} (no floor), and
\emph{control} (\texttt{White $\times$ Male} for HMDA, \texttt{40\_plus
$\times$ male} for German Credit). The body's Table~\ref{tab:cells} is the
first of these.

\subsection*{HMDA race $\times$ sex, reference = max rate (no floor)}
{\small\input{generated/appendix_a_hmda_max_rate}}

\subsection*{HMDA race $\times$ sex, reference = White $\times$ Male}
{\small\input{generated/appendix_a_hmda_control}}

\subsection*{German Credit AgeGroup $\times$ Sex, reference = max rate above floor}
{\small\input{generated/appendix_a_german_credit_max_rate_floor}}

\subsection*{German Credit AgeGroup $\times$ Sex, reference = max rate (no floor)}
{\small\input{generated/appendix_a_german_credit_max_rate}}

\subsection*{German Credit AgeGroup $\times$ Sex, reference = 40\_plus $\times$ male}
{\small\input{generated/appendix_a_german_credit_control}}

\medskip\noindent Marginal disparate impact, full population, for
completeness: German Credit \texttt{Sex} $0.8888$, \texttt{AgeGroup} $0.9258$,
\texttt{foreign\_worker} $0.8387$; HMDA \texttt{race} $0.9371$, \texttt{sex}
$0.9572$, \texttt{age\_group} (banded) $0.9791$, \texttt{age\_62\_plus}
$0.8897$; Lending Club (synthetic gender; instrument-test setting only)
\texttt{gender} $0.9931$, \texttt{income\_level} $1.0072$,
\texttt{loan\_amount\_level} $1.0063$. Source:
\texttt{artifacts/*/fairness/fairness\_summary.json}.

\section{Minimal reproduction of the favorable-label default}
\label{app:repro}

The script (deterministic: fixed counts, no RNG), its complete output, and
the pinned environment (\texttt{requirements-repro.txt}: \texttt{aif360==0.6.1},
\texttt{fairlearn==0.13.0}). Two groups of $1{,}772$; the privileged group
defaults $111$ times, the unprivileged $134$.

{\scriptsize\verbatiminput{../scripts/aif360_default_repro.py}}

\noindent Output (\texttt{artifacts/consolidated/aif360\_default\_repro.txt}):

{\scriptsize\verbatiminput{../artifacts/consolidated/aif360_default_repro.txt}}

\section{Subsampling study: design and full grid}
\label{app:subsample}

\paragraph{Design.} Population: the $10{,}978$ out-of-fold HMDA predictions
(random forest, stratified five-fold, seed 42). Draws: simple random, without
replacement, $B{=}2{,}000$ per size, seeded; note the historical $2{,}196$
split was outcome-stratified, so the study's draws are a slightly harder
regime than the original audit faced. Ground truth $\theta$ is the
full-population cell table \emph{under each reference convention}; violating
cells are those with $\theta$-DI below $0.80$ excluding the reference (four
cells under the max-rate conventions, three under the control convention).
Policies as in \S\ref{sec:subsample}. The Wilson policy's DI interval is
deliberately conservative (disadvantaged cell's lower bound over the
reference's upper bound, ignoring the correlation between the two estimates)
and it over-covers: measured coverage $0.993$--$1.0$ against nominal
$0.95$. That is the intended direction of error for an audit, and it is why
the abstention rates in the grid are upper bounds on what a sharper interval
(delta method, bootstrap) would produce \citep{chen2018}. The empirical-Bayes
policy shrinks cell rates toward the pooled rate with a beta prior whose mean
is the pooled rate and whose strength $\tau = (1-\rho)/\rho$ comes from the
weighted (Kleinman) beta-binomial moment estimator of the intraclass
correlation $\rho$, clipped to $[10^{-4}, 0.99]$; on this population
$\tau \approx 90$, which is why the policy clears small violating cells even
at full population strength. A data-dependent reference (both max-rate
conventions) is re-fit within every draw, because that is what the convention
\emph{is}, while the control reference is fixed.

\paragraph{Full grid.} Rates in \% of draws.

{\scriptsize\input{generated/subsample_full_grid}}

\section{Population construction}
\label{app:attrition}

Generated by \texttt{scripts/hmda\_attrition.py}, which replays
\texttt{clean\_hmda.py}'s row-dropping rules in order.

\subsection*{Attrition, 20{,}000 $\to$ 10{,}978}
{\small\input{generated/attrition_table}}

\subsection*{Race DI on raw approval rates, dropped categories retained}

Raw approvals (\texttt{action\_taken == 1}) over all $20{,}000$ rows: a
property of the data, not the model, since dropped rows never receive
predictions. The excluded \texttt{Race Not Available} group ($4{,}140$
applicants) approves at $0.7234$: it clears the $0.80$ screen against White
applicants ($0.9205$) and sits below every retained majority group, and its
exclusion is invisible in every downstream table of this paper.

{\small\input{generated/race_not_available_di}}

\section{The instruments, verbatim}
\label{app:instruments}

Everything in this appendix is extracted from the live code by
\texttt{scripts/emit\_appendix\_f.py} (via \texttt{inspect.getsource}) and
kept current by a fidelity test, so it cannot drift from what actually ran.
The frozen prompt packs (fully rendered, with data) are in
\texttt{artifacts/llm\_benchmark/dry\_run/}.

\subsection*{Cycle-4 decomposition backing \S\ref{sec:capability}}

Generated by \texttt{scripts/recheck\_prompt\_sensitivity.py}. ``Share of
dip'' columns decompose the constrained-cycle drop in mean total score by
subscore.

{\scriptsize\input{generated/prompt_sensitivity_decomposition}}

\input{generated/appendix_f_code}

\section{Claim $\to$ artifact $\to$ command}
\label{app:map}

Generated by \texttt{scripts/artifact\_map.py};
\texttt{tests/test\_artifact\_map.py} asserts every artifact exists and the
table is current. Two headline numbers were previously published with no
artifact behind them, and both were wrong; this table is the control that
prevents a third.

\begin{table*}[p]
\caption{Every load-bearing claim, the artifact it traces to, and the command
that regenerates it.}
\label{tab:map}
\scriptsize
% Generated by scripts/artifact_map.py -- do not edit.
\begin{tabular}{@{}p{5.2cm}p{5.4cm}p{5.6cm}@{}}
  \toprule
  claim & artifact & regenerate with \\
  \midrule
  AIF360 default DI 1.2072 vs oriented 0.9862; Fairlearn 0.8284 on the column as encoded & \texttt{artifacts/consolidated/aif360\_default\_repro.txt} & \texttt{python3.11 scripts/aif360\_default\_repro.py} \\
  \addlinespace
  In-vivo orientation artifact: DI 0.0 (CRITICAL) vs 0.9931 on Lending Club gender & \texttt{artifacts/lending\_club/fairness/fairness\_summary.json} & \texttt{python3.11 run\_pipeline.py german\_credit hmda lending\_club} \\
  \addlinespace
  Marginal DI per attribute (Table 2), incl. banded age 0.9791 vs age-62+ 0.8897 & \texttt{artifacts/hmda/fairness/fairness\_summary.json} & \texttt{python3.11 run\_pipeline.py german\_credit hmda lending\_club} \\
  \addlinespace
  Full-population race x sex cells with Wilson CIs; 6/12 indeterminate; defensible set & \texttt{artifacts/consolidated/interval\_estimates.json} & \texttt{python3.11 scripts/interval\_estimates.py --emit-tex} \\
  \addlinespace
  Reference-convention shift: 0.6734 to 0.7012, 0.8632 to 0.8988 vs White x Male & \texttt{artifacts/consolidated/interval\_estimates.json} & \texttt{python3.11 scripts/interval\_estimates.py --emit-tex} \\
  \addlinespace
  Split vs population: 5/11 cells suppressed; 1.1504 (n=4) to 0.6734 (n=31); 0.1917 to 0.8815 & \texttt{artifacts/consolidated/population\_comparison.json} & \texttt{python3.11 scripts/compare\_populations.py} \\
  \addlinespace
  Subsampling rates (Fig. 1, Table 3, subsampling appendix): false clearance 83.8\% point / 100.0\% EB; floors 100\% wrong worst cell; Wilson abstains 90.3\% (m=2196) & \texttt{artifacts/consolidated/subsample\_study.json} & \texttt{python3.11 scripts/subsample\_study.py} \\
  \addlinespace
  Population construction: 20,000 to 10,978; Race Not Available 4,140 rows (20.7\%); raw-approval DI with dropped categories retained & \texttt{artifacts/consolidated/population\_attrition.json} & \texttt{python3.11 scripts/hmda\_attrition.py} \\
  \addlinespace
  Refusal detector: 6 of 9 constrained-cycle responses flagged, zero refusals & \texttt{artifacts/llm\_benchmark/gpt-4o} & \texttt{python3.11 scripts/report\_track\_q.py} \\
  \addlinespace
  Gemini means 19.79 (all) / 47.50 (excl. zero scores) / 72.50 (content-bearing); 18 empty results; 4 empty-but-scored & \texttt{artifacts/consolidated/prompt\_sensitivity\_check.json} & \texttt{python3.11 scripts/recheck\_prompt\_sensitivity.py} \\
  \addlinespace
  Provider block recorded on the invalid measurement & \texttt{configs/research\_guardrails.json} & \texttt{cat configs/research\_guardrails.json} \\
  \addlinespace
  Fabrication flags 17 vs 4 on identical packs; 14 of 21 derivable from context & \texttt{artifacts/consolidated/RESULTS.md} & \texttt{python3.11 scripts/llm\_benchmark.py --replay (see RESULTS.md S4.1)} \\
  \addlinespace
  Track Q agreement 75.0\% decomposing to 33.3/91.7/100.0; 7 over- vs 2 under-escalations & \texttt{artifacts/consolidated/track\_q\_analysis.json} & \texttt{python3.11 scripts/report\_track\_q.py} \\
  \addlinespace
  Prompt-sensitivity spreads 14.67/9.0 vs 0.0/1.88; spread excl. constrained cycle 6.67/1.25; dip decomposition & \texttt{artifacts/consolidated/prompt\_sensitivity\_check.json} & \texttt{python3.11 scripts/recheck\_prompt\_sensitivity.py} \\
  \addlinespace
  Track H: 27/27 citations verified against the retrieved pool & \texttt{artifacts/consolidated/track\_h\_evaluation.json} & \texttt{python3.11 scripts/evaluate\_track\_h.py} \\
  \addlinespace
  German Credit AgeGroup x Sex cell table; under-40 female 0.8201 & \texttt{artifacts/consolidated/interval\_estimates.json} & \texttt{python3.11 scripts/interval\_estimates.py --emit-tex} \\
  \addlinespace
  Out-of-fold scoring: RF acc 0.8002 / AUC 0.8018 over 10,978 rows, 5-fold, seed 42 & \texttt{hmda\_dataset/metrics/out\_of\_fold\_metrics.json} & \texttt{python3.11 hmda\_dataset/scripts/train\_models.py --full-population} \\
  \addlinespace
  \bottomrule
\end{tabular}

\end{table*}

\end{document}
